\section{Introduction}
\label{sec:intro}

Text-to-video models such as Wan~\cite{wan2025wan}, CogVideoX~\cite{yang2025cogvideox} and VACE~\cite{jiang2025vace} synthesise convincing street scenes, shop fronts and screens, yet the text inside them is often garbled: a record shop reads \str{ROCRDS}, a convenience store reads \cjkimg{cjk_eab75b7f4d} (\cref{fig:teaser}). In our benchmark of generated clips, only \rOrigRead\% of the target strings are rendered correctly. For advertising, signage, education or any use in which the text carries meaning, a single wrong character is a hard failure, and re-generating the clip until the text happens to come out right is expensive and changes the shot. We therefore study \emph{repair}: keep the generated clip, and make its text read correctly.

The natural tool is a generative video text editor~\cite{zeng2026steervte}, which regenerates the text region conditioned on rendered glyphs and on the style of the source. Such editors produce glyphs that blend into the scene, but on generated videos we observe three systematic problems. They write characters that look plausible but are wrong---\str{FRESY BREAD}, \zh{小心坍滑}, \zh{洹影院}---and on our benchmark the strongest of them, \steer~\cite{zeng2026steervte}, still leaves \rSteerUnread\% of the clips unreadable. They rewrite text that was already correct, turning \str{FRESH BREAD} into \str{FRECH BREAD} (\cref{fig:teaser}c). And because the whole clip passes through the generator, the background changes as well: the PSNR outside the text region drops to \rSteerBG\,dB.

We argue that repair should be correctness-first and minimal, with a strict order of priorities: do not change what is already right; when a change is needed, change only the text; and recover the material of the old text only when there is evidence that this helps. \method implements this order with three evidence gates (\cref{sec:method}). A preservation gate returns the clip unchanged when OCR on many frames already reads the target. A localization gate decides whether the new text can be laid out on the old glyph instance---its size, colour and baseline---or must fall back to filling the tracked box. An appearance gate decides whether to apply source appearance transport (SAT), a non-parametric operator that carries glow and material from the old glyphs to the new ones. Between the gates, the target string is rendered deterministically from a matched font and warped with carrier homographies, so its spelling cannot drift, and only pixels inside the text region are composited. The pipeline requires no training.

On 526 generated clips covering seven carrier materials in English and Chinese, \method raises readability from \rOrigRead\% to \rOursRead\% (\steer: \rSteerRead\%), with \rOursNet{} net repairs against \rSteerNet, and keeps the background intact (\rOursBG\,dB). We are equally explicit about what it does not do. Its glyphs are rendered, not generated: where both methods repair a garbled clip, \steer's text is closer to the original appearance in \rSteerCloserPct\% of the clips, and \method breaks \rOursBroken{} originally correct clips where \steer breaks \rSteerBroken. We trace these gaps to their causes (\cref{sec:limits}) and report the alternatives we tried and rejected, including two modules aimed directly at the appearance gap that did not pass their pre-registered criteria (\cref{sec:exp_neg}).

Our contributions are:
\begin{itemize}
\item \textbf{A repair task, benchmark and protocol.} 526 generated clips with hint boxes and targets, seven materials, two languages; metrics that separate repairs from breakages (net repair), background fidelity and collateral change, temporal stability, and blind human judgements. Development clips are disclosed and excluded in a separate report; the selective-editing study, the two appearance modules of \cref{sec:exp_neg} and the second-reader evaluation were pre-registered.
\item \textbf{\method}, a training-free, evidence-gated repair pipeline: preservation and localization gates, deterministic glyph rendering anchored to the tracked carrier, and SAT behind an appearance gate.
\item \textbf{An extensive analysis including negative results}: a learned 1.3B refiner trained on up to one million pairs does not improve appearance and costs readability and temporal stability; glyph-only erasure and a carrier-fit check fail their pre-registered criteria; reuse of the correctly generated source glyphs changes too few clips with our current segmentation, matching and gating to form a general module, baseline-aligned layout with glyph-matched fonts shows no overall benefit under our evaluation protocol, whose appearance metric is blind to glyph shape, so its perceptual effect remains unverified; and a glyph-locked material generator trained on rendered counterfactual pairs keeps the glyphs but, under our registered protocol, does not transfer zero-shot to generated video; and a failure taxonomy that locates the remaining appearance gap. Because the two systems fail on different clips, we also register and evaluate a verified selection between them: a reader independent of the scorer accepts \steer's text where it reads correctly and composites it locally into the original, which reaches \hVRead\% readability with \hVBroken{} broken clips (\cref{sec:exp_hybrid}).
\end{itemize}
